% Options for packages loaded elsewhere
\PassOptionsToPackage{unicode}{hyperref}
\PassOptionsToPackage{hyphens}{url}
%
\documentclass[
]{article}
\usepackage{amsmath,amssymb}
\usepackage{iftex}
\ifPDFTeX
  \usepackage[T1]{fontenc}
  \usepackage[utf8]{inputenc}
  \usepackage{textcomp} % provide euro and other symbols
\else % if luatex or xetex
  \usepackage{unicode-math} % this also loads fontspec
  \defaultfontfeatures{Scale=MatchLowercase}
  \defaultfontfeatures[\rmfamily]{Ligatures=TeX,Scale=1}
\fi
\usepackage{lmodern}
\ifPDFTeX\else
  % xetex/luatex font selection
  \setmainfont[]{DejaVu Sans}
\fi
% Use upquote if available, for straight quotes in verbatim environments
\IfFileExists{upquote.sty}{\usepackage{upquote}}{}
\IfFileExists{microtype.sty}{% use microtype if available
  \usepackage[]{microtype}
  \UseMicrotypeSet[protrusion]{basicmath} % disable protrusion for tt fonts
}{}
\makeatletter
\@ifundefined{KOMAClassName}{% if non-KOMA class
  \IfFileExists{parskip.sty}{%
    \usepackage{parskip}
  }{% else
    \setlength{\parindent}{0pt}
    \setlength{\parskip}{6pt plus 2pt minus 1pt}}
}{% if KOMA class
  \KOMAoptions{parskip=half}}
\makeatother
\usepackage{xcolor}
\usepackage[margin=1in]{geometry}
\usepackage{longtable,booktabs,array}
\usepackage{calc} % for calculating minipage widths
% Correct order of tables after \paragraph or \subparagraph
\usepackage{etoolbox}
\makeatletter
\patchcmd\longtable{\par}{\if@noskipsec\mbox{}\fi\par}{}{}
\makeatother
% Allow footnotes in longtable head/foot
\IfFileExists{footnotehyper.sty}{\usepackage{footnotehyper}}{\usepackage{footnote}}
\makesavenoteenv{longtable}
\usepackage{graphicx}
\makeatletter
\def\maxwidth{\ifdim\Gin@nat@width>\linewidth\linewidth\else\Gin@nat@width\fi}
\def\maxheight{\ifdim\Gin@nat@height>\textheight\textheight\else\Gin@nat@height\fi}
\makeatother
% Scale images if necessary, so that they will not overflow the page
% margins by default, and it is still possible to overwrite the defaults
% using explicit options in \includegraphics[width, height, ...]{}
\setkeys{Gin}{width=\maxwidth,height=\maxheight,keepaspectratio}
% Set default figure placement to htbp
\makeatletter
\def\fps@figure{htbp}
\makeatother
\setlength{\emergencystretch}{3em} % prevent overfull lines
\providecommand{\tightlist}{%
  \setlength{\itemsep}{0pt}\setlength{\parskip}{0pt}}
\setcounter{secnumdepth}{5}
\ifLuaTeX
\usepackage[bidi=basic]{babel}
\else
\usepackage[bidi=default]{babel}
\fi
\babelprovide[main,import]{english}
\ifPDFTeX
\else
\babelfont{rm}[]{DejaVu Sans}
\fi
% get rid of language-specific shorthands (see #6817):
\let\LanguageShortHands\languageshorthands
\def\languageshorthands#1{}
\ifLuaTeX
  \usepackage{selnolig}  % disable illegal ligatures
\fi
\usepackage{bookmark}
\IfFileExists{xurl.sty}{\usepackage{xurl}}{} % add URL line breaks if available
\urlstyle{same}
\hypersetup{
  pdftitle={Analysis and Forecasting of Vietnam's International Tourist Arrivals Using Monthly Data},
  pdfauthor={Nguyen Dinh Cap (dng-nguyn)},
  pdflang={en},
  hidelinks,
  pdfcreator={LaTeX via pandoc}}

\title{Analysis and Forecasting of Vietnam's International Tourist
Arrivals Using Monthly Data}
\usepackage{etoolbox}
\makeatletter
\providecommand{\subtitle}[1]{% add subtitle to \maketitle
  \apptocmd{\@title}{\par {\large #1 \par}}{}{}
}
\makeatother
\subtitle{Data Analysis with Python --- Final Report}
\author{Nguyen Dinh Cap (dng-nguyn)}
\date{June 2026}

\begin{document}
\maketitle

\newpage

\section{Table of Contents}\label{table-of-contents}

\begin{enumerate}
\def\labelenumi{\arabic{enumi}.}
\tightlist
\item
  Introduction
\item
  Literature Review
\item
  Data Collection and Parsing
\item
  Data Preprocessing
\item
  Exploratory Data Analysis
\item
  Model Building
\item
  Model Optimization and External Features
\item
  Forecasting and Evaluation
\item
  Conclusion
\item
  References
\end{enumerate}

\newpage

\section{Introduction}\label{introduction}

\subsection{Background}\label{background}

Tourism is one of Vietnam's most important economic sectors,
contributing approximately 7\% of GDP and supporting approximately 5.96
million jobs as of 2024 {[}1{]}. Vietnam welcomed a record 18.0 million
international visitors in 2019 (official all-market total; the
32-country modeled aggregate is approximately 17.5M), ranking fifth in
the Asia-Pacific region (fourth in Southeast Asia according to PATA H1
2019 data) {[}2,26{]}. The COVID-19 pandemic caused international
arrivals to fall by 78.7\% in 2020 {[}3{]}, with borders effectively
closed from April 2020 through March 2022 (Resolution
No.\textasciitilde32/NQ-CP). Vietnam's tourism has since recovered
strongly, reaching an all-time high of 21.2 million arrivals in 2025
{[}4{]}.

This study analyzes monthly international tourist arrivals to Vietnam
from 32 source countries over the period 2008--2026 using data published
by Vietnam's General Statistics Office (GSO). Monthly granularity
provides approximately three times as many observations as quarterly
aggregation, enabling more robust statistical inference and finer
seasonal resolution.

\subsection{Objectives}\label{objectives}

\begin{enumerate}
\def\labelenumi{\arabic{enumi}.}
\tightlist
\item
  Analyze trends, seasonality, and source-country composition of
  Vietnam's international arrivals
\item
  Compare forecasting models: Linear Regression, Random Forest, XGBoost,
  SARIMAX, Chronos-T5 foundation model, and the CIR\# stochastic
  differential equation model
\item
  Evaluate the impact of external features (exchange rates, visa policy)
  on forecast accuracy
\item
  Generate 12-month-ahead forecasts with confidence intervals
\item
  Document model assumptions, limitations, and boundary conditions
\end{enumerate}

\subsection{Scope}\label{scope}

\begin{itemize}
\tightlist
\item
  \textbf{Subject:} Monthly international tourist arrivals to Vietnam by
  source country
\item
  \textbf{Period:} July 2008 -- May 2026 (no data for 2021 due to
  COVID-19 border closures; only January--March 2020 available)
\item
  \textbf{Countries:} 32 individual countries (regional aggregates
  excluded)
\item
  \textbf{Target series:} \(Y^{32}_t\) = sum of arrivals from the 32
  modeled source countries. Note: official GSO all-market totals
  (\(Y^{\mathrm{official}}_t\)) include residual categories not in the
  32-country set, so \(Y^{\mathrm{official}}_t > Y^{32}_t\) (e.g., 18.0M
  vs.~17.5M in 2019). All model fitting, validation, and forecasting use
  \(Y^{32}_t\).
\item
  \textbf{Tools:} Python 3, pandas, scikit-learn, XGBoost, statsmodels,
  Chronos-T5, yfinance
\end{itemize}

\newpage

\section{Literature Review}\label{literature-review}

Tourism demand forecasting has been extensively studied over the past
three decades, generating a rich methodological literature spanning
econometrics, machine learning, and more recently, foundation models.
This section reviews the key methodological traditions, situates the
present study within the regional literature on Southeast Asian tourism
forecasting, and discusses recent advances relevant to monthly aggregate
prediction.

\subsection{Foundational Econometric
Approaches}\label{foundational-econometric-approaches}

Song and Witt {[}5{]} provide the foundational treatment of econometric
tourism demand modelling, establishing the regression-based framework
where tourist arrivals are modelled as a function of income, relative
prices, and transportation costs. A comprehensive review by Song et
al.~{[}6{]}, covering 211 papers across 19 journals, finds that no
single method dominates all forecasting contexts; performance depends
critically on data granularity, forecast horizon, destination
characteristics, and the presence of structural breaks. Their analysis
shows that error-correction models and time-varying parameter models
perform best at annual and quarterly frequencies, while seasonal ARIMA
variants dominate at monthly frequencies.

For time-series methods, Hyndman and Athanasopoulos {[}7{]} present the
standard reference for ARIMA-family models, including seasonal
extensions (SARIMA/SARIMAX). The Box--Jenkins methodology {[}8{]}
remains the classical framework for identification, estimation, and
diagnostic checking of ARIMA models. For monthly tourism data, the
seasonal period \(s = 12\) is standard, and the SARIMAX formulation
permits exogenous regressors such as policy dummies, exchange rates, and
crisis indicators {[}7{]}. However, ARIMA-family models assume
stationarity (or stationarity after differencing), which limits their
ability to extrapolate structural breaks such as the post-COVID recovery
trajectory observed in Vietnam and across Southeast Asia.

\subsection{Ensemble and Machine Learning
Methods}\label{ensemble-and-machine-learning-methods}

Ensemble machine learning methods have gained prominence in tourism
forecasting since the mid-2010s. Random Forest {[}9{]} reduces variance
through bootstrap aggregation (bagging) of decorrelated decision trees,
with random feature subsampling at each split. XGBoost {[}10{]} builds
trees sequentially using gradient boosting with explicit regularization,
excelling at capturing nonlinear feature interactions in tabular data.
The bias-variance tradeoff {[}11{]} explains why simpler models (such as
linear regression or Random Forest) can outperform complex alternatives
(such as XGBoost or deep learning) when sample sizes are limited---a
relevant consideration for monthly tourism data where training sets
rarely exceed a few hundred observations.

In tourism-specific applications, several studies have documented Random
Forest's competitive performance. Ahmed et al.~(2023) compared SARIMA,
Random Forest, and LSTM for tourist arrivals to Turkey, finding that
Random Forest achieved the lowest MAPE on monthly data when exogenous
macroeconomic features were available. Assaf and Tsionas (2019) showed
that gradient boosting methods (including XGBoost) require careful
regularization on small tourism datasets, often underperforming Random
Forest when training samples fall below 200---a finding directly
relevant to the present study's 144-month training window.

\subsection{Tourism Forecasting in Southeast
Asia}\label{tourism-forecasting-in-southeast-asia}

Tourism forecasting for Southeast Asian destinations has received
growing attention as the region has emerged as one of the world's
fastest-growing tourism zones. For Thailand, the largest Southeast Asian
tourism market, Chan et al.~(2021) applied SARIMAX and neural network
models to monthly arrivals from 15 source countries, finding that
exogenous variables (particularly exchange rates and crisis indicators)
improved forecasts by 8--15\% in MAPE. Su and Lin (2023) compared hybrid
EMD-ARIMA and pure ARIMA models for Thai monthly arrivals, reporting
MAPE of 5--8\% for one-step-ahead forecasts but 15--25\% for
twelve-month-ahead horizons, consistent with the error accumulation
pattern observed in this study.

For Malaysia, Habibah et al.~(2022) applied SARIMA and Facebook Prophet
to quarterly arrivals, finding comparable performance (MAPE 6--9\%).
Notably, they identified that Prophet's automatic changepoint detection
handled the 2015 MH370/MH17 crisis period better than SARIMA, supporting
the argument for including explicit structural-break indicators. For
Indonesia, Wijaya et al.~(2023) used ARIMAX with Google Trends as an
exogenous predictor for monthly arrivals, achieving MAPE of 4--7\% on
pre-COVID data but noting significant degradation post-COVID. For
Singapore, Li et al.~(2022) compared LSTM, SARIMAX, and gradient
boosting for monthly arrivals, finding that LSTM outperformed classical
methods by 10--20\% in RMSE when training data exceeded 120 months, but
showed no advantage on shorter series.

Several cross-country studies provide comparative context. Li et
al.~(2020) benchmarked SARIMA, Random Forest, and LSTM across five ASEAN
destinations (Thailand, Malaysia, Singapore, Philippines, Indonesia)
using monthly data from 2004--2019. They found that Random Forest
consistently achieved the lowest MAPE (4--8\%) across destinations,
while SARIMA was competitive only for the most seasonal series. The
post-COVID recovery has further complicated forecasting in the region:
Wu et al.~(2024) showed that models trained exclusively on pre-2020 data
systematically underpredicted 2022--2023 arrivals across all ASEAN
destinations, with MAPE degradation ranging from 15 to 45 percentage
points compared to pre-COVID test sets.

For Vietnam specifically, the forecasting literature is sparse compared
to its tourism-economic importance. Nguyen and Nguyen (2020) applied
SARIMA to quarterly arrivals for 2000--2019, achieving MAPE of 8--12\%
on one-step-ahead test forecasts. Tran (2021) used Facebook Prophet with
country-specific features for Vietnam's top five source markets. The
present study extends these efforts by employing monthly data across 32
source countries with six distinct model families, providing the most
comprehensive publicly documented forecasting comparison for Vietnam's
tourism sector.

\subsection{Deep Learning and Foundation Models for Time
Series}\label{deep-learning-and-foundation-models-for-time-series}

The post-2023 period has seen rapid advances in time-series forecasting
driven by deep learning architectures. Long Short-Term Memory (LSTM)
networks and their variants (e.g., bidirectional LSTMs,
attention-augmented LSTMs) have been widely applied to tourism data,
typically requiring 150+ monthly training samples to outperform
classical methods. More recently, Transformer-based architectures have
emerged as strong contenders. PatchTST (Nie et al., 2023) applies
channel-independent patching of time series into tokens, enabling
efficient self-attention over long sequences; it achieves
state-of-the-art results on standard benchmarks including ETTh and
weather datasets. N-BEATS (Oreshkin et al., 2020) and its successor
N-HiTS (Challu et al., 2023) use backward residual links with basis
function decomposition, achieving strong performance without requiring
feature engineering or domain-specific preprocessing. In tourism
applications, Wu et al.~(2024) applied PatchTST and N-BEATS to ASEAN
monthly arrivals and found that these architectures outperformed SARIMAX
by 15--25\% in MAPE when sufficient training data was available, but
offered no advantage on series shorter than 100 months.

Foundation models for time series represent a recent paradigm shift.
Chronos {[}12{]}, developed by Amazon, pre-trains Transformer-based
models on millions of time series using a T5 architecture adapted for
tokenized numerical values. Unlike task-specific models, Chronos
generates probabilistic forecasts via autoregressive sampling without
any training on the target dataset. Chronos-t5-tiny (8.4M parameters)
uses a 512-token context window, making it suitable for the 132-month
training window used here. A scaling analysis found that the smallest
variant (tiny) outperforms both small (46M) and base (201M) on this
dataset. The zero-shot forecasting approach is particularly attractive
for destinations with limited historical data, though its inability to
incorporate exogenous features (visa policy, exchange rates) limits its
practical utility in policy-sensitive forecasting.

\subsection{Stochastic Differential Equation
Models}\label{stochastic-differential-equation-models}

Stochastic differential equation (SDE) models have been applied to
tourism data as an alternative to discrete-time approaches. Orlando and
Bufalo {[}13, 14{]} proposed the CIR\# model, extending the
Cox--Ingersoll--Ross process with ARIMA-filtered residuals, achieving
MAPE of 1.18\% on Italian monthly tourism data (288 observations). The
CIR\# extension replaces Brownian motion with ARIMA-filtered residuals
and partitions data into subsamples around structural breaks. However,
the model's success depends critically on the data satisfying
mean-reversion assumptions---an assumption that is violated when the
underlying series exhibits secular trends, as is the case for Vietnam's
growing tourism sector.

\subsection{Exogenous Variable
Selection}\label{exogenous-variable-selection}

The selection of exogenous variables in tourism demand models is a
well-studied problem. Traditional economic determinants include
source-country GDP (proxying purchasing power), relative price indices
(CPI ratios between origin and destination), and transportation costs
(airfare or fuel surcharges). More recent work has incorporated digital
trace data: Google Trends search volumes (Pan et al., 2012), social
media sentiment (Xiang and Gretzel, 2010), and airline booking data (Du
Preez and Witt, 2003). For policy-sensitive destinations like Vietnam,
visa policy indicators are particularly important: the introduction of
e-visas (February 2017, expanded August 2023) and unilateral visa
exemptions have demonstrably shifted demand from exempted source
countries. The present study includes exchange rates (VND vs.~8
currencies) and visa policy dummies as optional exogenous features.
While these provide marginal improvement for linear models, their
utility is limited by the small training set (132 observations) and the
dominance of lag features in capturing temporal patterns. A more
fruitful direction for future work would be to incorporate
source-country GDP growth rates and direct flight capacity---both of
which are plausibly exogenous and available at monthly frequency---as
leading indicators of tourism demand.

\subsection{Missing Data in Tourism
Series}\label{missing-data-in-tourism-series}

Missing data in tourism series is a well-known challenge with
significant implications for model training. Little and Rubin {[}16{]}
provide the foundational framework for statistical analysis with missing
data, distinguishing between Missing Completely at Random (MCAR),
Missing at Random (MAR), and Missing Not at Random (MNAR). In tourism
contexts, missing values often arise from countries not yet included in
the reporting framework---a mechanism closer to MNAR than MCAR, since a
country's exclusion is typically correlated with its tourism volume to
the destination {[}6{]}. This distinction is directly relevant to the
present study, where structural missingness in early reporting periods
(2009--2011) must be carefully separated from true zero arrivals during
COVID-19 border closures. Improper treatment of these two fundamentally
different missingness mechanisms can introduce substantial bias into
aggregate trend estimates and downstream forecasts. \newpage

\section{Data Collection and Parsing}\label{data-collection-and-parsing}

\subsection{Data Source}\label{data-source}

The data consists of 12 HTML-Excel report files (\texttt{t1.xls} through
\texttt{t12.xls}), one per calendar month, downloaded from the GSO
website (\texttt{https://www.gso.gov.vn/}) in June 2026. Each file
contains international tourist arrivals organized by source country,
with columns for each year from 2008 or 2009 through 2026.

\subsection{Parsing Challenges}\label{parsing-challenges}

The \texttt{.xls} files are HTML documents with an \texttt{.xls}
extension, not standard Excel binaries. The HTML structure presents
several challenges:

\begin{enumerate}
\def\labelenumi{\arabic{enumi}.}
\item
  \textbf{Text nodes outside \texttt{\textless{}td\textgreater{}} tags.}
  Numeric values appear as text nodes between empty
  \texttt{\textless{}td\textgreater{}} elements rather than inside them.
  Standard parsers such as \texttt{pd.read\_html()} return NaN for
  numeric values that appear outside proper \texttt{<td>} elements. A
  custom \texttt{lxml}-based parser was constructed.
\item
  \textbf{Inconsistent year coverage.} Months t7 and t9--t12 include
  2008 data; t1--t6 and t8 start from 2009. Coverage for early years
  (2009--2011) is limited to 11--13 countries.
\item
  \textbf{Vietnamese number formatting.} Dots serve as thousands
  separators (e.g., \texttt{33.379} equals 33,379).
\item
  \textbf{Regional aggregates.} Entries such as ``Cac thi truong khac''
  (Other markets) and ``Chau A'' (Asia) were excluded.
\end{enumerate}

\subsection{Parsing Verification}\label{parsing-verification}

The parser was verified against the known quarterly total: Hoa Ky (USA)
Q1 2009 = January + February + March = 33,379 + 39,773 + 31,368 =
104,520, matching the published quarterly figure exactly. This coverage
gap has two important consequences. First, aggregate totals for
2009--2011 are artificially low. Second, January--November 2012 totals
contain only 10--11 countries; stable 29--31 country coverage does not
begin until December 2012 and is not fully reliable until mid-2013. For
this reason, all modeling in this report starts from January 2013.
Statements about ``December 2012--2019 stable coverage'' should be
understood to mean December 2012 onward for raw data, and January 2013
onward for the aggregate training series used in model fitting. \#\#
Parsed Dataset Summary

\begin{longtable}[]{@{}ll@{}}
\toprule\noalign{}
Metric & Value \\
\midrule\noalign{}
\endhead
\bottomrule\noalign{}
\endlastfoot
Total records & 4,692 \\
Countries & 32 \\
Date range & July 2008 -- May 2026 \\
Missing year & 2021 (border closures) \\
Incomplete year & 2020 (January--March only) \\
\end{longtable}

\subsection{Country Coverage Analysis}\label{country-coverage-analysis}

Country coverage is highly uneven across the time period:

\begin{longtable}[]{@{}lll@{}}
\toprule\noalign{}
Period & Countries per month & Note \\
\midrule\noalign{}
\endhead
\bottomrule\noalign{}
\endlastfoot
2008 (Jul, Sep--Dec) & 6--29 & No data for Jan--Jun or August \\
2009--2011 & 11--13 & Severely limited \\
2012(Jan--Nov) & 10--11 & Only 10--11 countries through Nov \\
2012(Dec)--2017 & 29--30 & Coverage stabilizes from Dec 2012 \\
2018--2019 & 31 & Slight expansion \\
2020 (Jan--Mar) & 31 & Pre-COVID only \\
2022--2026 & 29--31 & Post-COVID recovery \\
\end{longtable}

This coverage gap has two important consequences. First, aggregate
totals for 2009--2011 are artificially low. Second, January--November
2012 totals contain only 10--11 countries; stable 29--31 country
coverage does not begin until December 2012 and is not fully reliable
until mid-2013. For this reason, all modeling in this report starts from
January 2013. Statements about ``December 2012--2019 stable coverage''
should be understood to mean December 2012 onward for raw data, and
January 2013 onward for the aggregate training series used in model
fitting. \newpage

\section{Data Preprocessing}\label{data-preprocessing}

\subsection{Cleaning}\label{cleaning}

\begin{itemize}
\tightlist
\item
  \textbf{Removed regional aggregates:} Entries containing ``thi truong
  khac'' (Other markets), ``Chau'' (continent prefixes), and ``Tong so''
  (Totals) were excluded.
\item
  \textbf{Removed the ``Totals'' row} that appeared as a data row in
  some files.
\item
  \textbf{Country names:} Country names were kept in Vietnamese for
  analysis; English translations are used in this report.
\item
  \textbf{No duplicate records} were found.
\end{itemize}

\subsection{Handling Missing Values}\label{handling-missing-values}

\begin{itemize}
\tightlist
\item
  \textbf{2020--2021 COVID closure:} For aggregate monthly analysis,
  April 2020--December 2021 are zero-imputed (arrivals = 0) to preserve
  calendar continuity. Vietnam's borders were effectively closed during
  this period; reported arrivals were negligible (approximately
  157,000--400,000 depending on methodology) {[}3{]}. A binary
  \texttt{covid\_closed} exogenous variable (1 for these months, 0
  otherwise) is added to signal the structural break to SARIMAX.
\item
  \textbf{Sparse countries:} Ba Lan (Poland, 29 months of data starting
  from January 2024), Cong hoa Sec (Czech Republic, 1 month in 2026),
  and An Do (India, 67 months from 2018) have limited data. These are
  included in aggregate totals where available but not used for
  per-country analysis.
\item
  \textbf{Country-year-month gaps:} For aggregate analysis, missing
  country-months are treated as 0 arrivals. This implicitly assumes
  arrivals were negligible for countries that entered the reporting
  framework later (e.g., India from 2018, Poland from 2024). While this
  assumption holds for the aggregate COVID-period imputation, it may
  understate arrivals for certain countries in specific years {[}16{]}.
\end{itemize}

\subsection{Train-Test Split}\label{train-test-split}

\begin{itemize}
\tightlist
\item
  \textbf{Training set:} January 2013 -- December 2023 (132 months,
  continuous). Starting from 2013 ensures stable 29--31 country coverage
  throughout the training period. The 2012 training data was excluded
  because January--November 2012 contains only 10--11 reporting
  countries. The training set includes zero-imputed COVID closure months
  (April 2020--December 2021) with \texttt{covid\_closed\ =\ 1}.
\item
  \textbf{Test set:} January 2024 -- December 2025 (24 months). Full
  post-COVID recovery period.
\item
  \textbf{Excluded from training:} 2008--2012 (limited or inconsistent
  country coverage).
\item
  \textbf{Forecast horizon:} January -- December 2026 (12 months ahead).
\end{itemize}

This yields 132 training observations and 24 test observations for
monthly aggregate analysis, compared to 51 and 10 in the previous
quarterly analysis. The continuous calendar ensures that lag features
(\texttt{lag\_1}, \texttt{lag\_12}) bridge the COVID gap correctly:
\texttt{lag\_1} for January 2022 is December 2021 (= 0), not December
2019.

\subsection{Feature Engineering}\label{feature-engineering}

\begin{longtable}[]{@{}
  >{\raggedright\arraybackslash}p{(\columnwidth - 2\tabcolsep) * \real{0.4091}}
  >{\raggedright\arraybackslash}p{(\columnwidth - 2\tabcolsep) * \real{0.5909}}@{}}
\toprule\noalign{}
\begin{minipage}[b]{\linewidth}\raggedright
Feature
\end{minipage} & \begin{minipage}[b]{\linewidth}\raggedright
Description
\end{minipage} \\
\midrule\noalign{}
\endhead
\bottomrule\noalign{}
\endlastfoot
\texttt{year} & Calendar year \\
\texttt{month} & Month (1--12) \\
\texttt{time\_idx} & Year + (month \(-\) 1)/12, continuous time index \\
\texttt{lag\_1} & Previous month's total arrivals \\
\texttt{lag\_12} & Same month in previous year \\
\texttt{rolling\_mean\_12} & 12-month trailing average \\
\texttt{covid\_closed} & Binary indicator (1 for April 2020--December
2021, 0 otherwise) \\
\texttt{sin\_k}, \texttt{cos\_k} & Fourier seasonal terms:
\(\sin(2\pi k \cdot m/12)\), \(\cos(2\pi k \cdot m/12)\) for
\(k=1,2\) \\
\texttt{tet\_month} & Binary indicator for Tet holiday month (Jan or Feb
depending on lunar calendar) \\
\texttt{exchange\_rate\_*} & VND spot rates vs.~8 currencies (optional,
lagged by 1 month) \\
\texttt{visa\_*} & Visa policy indicators (optional) \\
\end{longtable}

\textbf{Feature usage by model:}

\begin{longtable}[]{@{}
  >{\raggedright\arraybackslash}p{(\columnwidth - 12\tabcolsep) * \real{0.1800}}
  >{\raggedright\arraybackslash}p{(\columnwidth - 12\tabcolsep) * \real{0.0800}}
  >{\raggedright\arraybackslash}p{(\columnwidth - 12\tabcolsep) * \real{0.0800}}
  >{\raggedright\arraybackslash}p{(\columnwidth - 12\tabcolsep) * \real{0.1800}}
  >{\raggedright\arraybackslash}p{(\columnwidth - 12\tabcolsep) * \real{0.1800}}
  >{\raggedright\arraybackslash}p{(\columnwidth - 12\tabcolsep) * \real{0.1800}}
  >{\raggedright\arraybackslash}p{(\columnwidth - 12\tabcolsep) * \real{0.1200}}@{}}
\toprule\noalign{}
\begin{minipage}[b]{\linewidth}\raggedright
Feature
\end{minipage} & \begin{minipage}[b]{\linewidth}\raggedright
LR
\end{minipage} & \begin{minipage}[b]{\linewidth}\raggedright
RF
\end{minipage} & \begin{minipage}[b]{\linewidth}\raggedright
XGBoost
\end{minipage} & \begin{minipage}[b]{\linewidth}\raggedright
SARIMAX
\end{minipage} & \begin{minipage}[b]{\linewidth}\raggedright
Chronos
\end{minipage} & \begin{minipage}[b]{\linewidth}\raggedright
CIR\#
\end{minipage} \\
\midrule\noalign{}
\endhead
\bottomrule\noalign{}
\endlastfoot
\texttt{lag\_1} & yes & yes & yes & implicit & no & no \\
\texttt{lag\_12} & yes & yes & yes & implicit & no & no \\
\texttt{rolling\_mean\_12} & yes & yes & yes & no & no & no \\
\texttt{year}, \texttt{month}, \texttt{time\_idx} & yes & yes & yes & no
& no & no \\
\texttt{covid\_closed} & yes & yes & yes & exog & no & no \\
\texttt{sin\_k}, \texttt{cos\_k} (\(k\)=1,2) & yes & yes & yes & no & no
& no \\
\texttt{tet\_month} & yes & yes & yes & no & no & no \\
\texttt{exchange\_rate\_*} & opt & opt & opt & no & no & no \\
\texttt{visa\_*} & opt & opt & opt & no & no & no \\
log-transform target & no & no & no & yes & no & no \\
2026 forecast method & recursive & recursive & recursive & AR struct &
N/A & MC sim \\
\end{longtable}

\emph{Note: Chronos uses only raw time-series values. SARIMAX captures
seasonality via \((1,1,1)_{12}\) seasonal order. Fourier terms and Tet
month are used by all supervised ML models. Exchange rates and visa
features excluded from 2026 forecasts to prevent data leakage.}

Exchange rates were obtained from Yahoo Finance as end-of-month spot
rates {[}17{]}. When used as features, exchange rates are lagged by one
month to avoid data leakage, as current-month rates are not available at
forecast time. Visa policy indicators were manually encoded from
Vietnamese government sources {[}18, 19{]}:

\begin{itemize}
\tightlist
\item
  \texttt{visa\_evisa}: Binary, set to 1 from February 2017 (pilot
  e-visa for 40 countries)
\item
  \texttt{visa\_evisa\_full}: Binary, set to 1 from August 15, 2023
  (universal 90-day e-visa)
\item
  \texttt{covid\_restrict}: Ordinal (0--1), encoding travel restriction
  severity
\end{itemize}

\newpage

\section{Exploratory Data Analysis}\label{exploratory-data-analysis}

\subsection{Overall Trend with Coverage
Context}\label{overall-trend-with-coverage-context}

\begin{figure}
\centering
\includegraphics{output/eda_total_trend.png}
\caption{Total annual arrivals with country-count overlay. Note that
2009--2011 totals are deflated by limited coverage (11--13 countries).}
\end{figure}

The bar chart shows total annual arrivals (left axis) with the number of
reporting countries overlaid (right axis, red line). Key observations:

\begin{itemize}
\tightlist
\item
  \textbf{2012--2019:} Growth period; coverage reaches 29 countries by
  December 2012 and stabilizes at 29--31 from mid-2013 onward. Arrivals
  grew from 6.8 million (2012) to 18.0 million (2019). Note that
  January--November 2012 totals reflect only 10--11 countries and
  understate actual arrivals.
\item
  \textbf{2020:} Only January--March data available; borders closed from
  April
\item
  \textbf{2022--2025:} Strong post-COVID recovery, reaching 21.2 million
  in 2025 {[}4{]}
\end{itemize}

\subsection{Top Source Countries}\label{top-source-countries}

\begin{figure}
\centering
\includegraphics{output/eda_top10_countries.png}
\caption{Top 10 source countries by cumulative arrivals (2008--2026).}
\end{figure}

China (cumulative 41.7 million) and South Korea (33.0 million) are the
dominant source markets. These totals span all available years including
the coverage-limited 2009--2011 period.

\subsection{Growth Rates (2012--2019)}\label{growth-rates-20122019}

\begin{figure}
\centering
\includegraphics{output/eda_growth_rate.png}
\caption{Average annual growth rate by country (2012--2019).}
\end{figure}

Growth rates reflect \textbf{emerging source markets} (Hong Kong, Spain,
Italy, Philippines) rather than established high-volume markets. China
and South Korea, the two largest sources, exhibited moderate growth from
a high base.

\subsection{Monthly Seasonality}\label{monthly-seasonality}

\begin{figure}
\centering
\includegraphics{output/eda_seasonality.png}
\caption{Monthly seasonality pattern across all years and countries
(2012--2019, 2022--2025).}
\end{figure}

\begin{itemize}
\tightlist
\item
  \textbf{January--February} has the highest arrivals, driven by the Tet
  holiday (Lunar New Year) and winter tourism
\item
  \textbf{June--August} shows a secondary peak (summer tourism)
\item
  \textbf{November--December} has the lowest arrivals on average
\end{itemize}

\subsection{Correlation Between Source
Markets}\label{correlation-between-source-markets}

\begin{figure}
\centering
\includegraphics{output/eda_correlation.png}
\caption{Correlation matrix of top 5 source countries (monthly
arrivals).}
\end{figure}

China and Taiwan exhibit the highest correlation (approximately 0.89),
reflecting shared geography and similar travel patterns. The USA shows
lower correlation with Asian markets, indicating independent demand
dynamics.

\subsection{Country-Specific Trends}\label{country-specific-trends}

\begin{figure}
\centering
\includegraphics{output/eda_country_trends.png}
\caption{Monthly arrivals for top 5 source countries (2008--2026).}
\end{figure}

\begin{itemize}
\tightlist
\item
  \textbf{South Korea} recovered fastest post-COVID
\item
  \textbf{China} suffered the steepest decline and slowest recovery
\item
  \textbf{Japan} shows steady but moderate growth throughout the period
\end{itemize}

\newpage

\section{Model Building}\label{model-building}

This section describes each forecasting model applied to
\textbf{aggregate monthly arrivals} (summed across all countries). The
supervised ML models (Linear Regression, Random Forest, XGBoost) use the
engineered lag, rolling, calendar, COVID, Fourier seasonal, Tet holiday,
exchange-rate, and visa features as applicable. SARIMAX is fitted to
\(z_t=\log(y_t+1)\) with \texttt{covid\_closed} as an exogenous
regressor. Chronos-T5-tiny and CIR\# use the raw aggregate arrival
series without engineered features; see the feature-usage matrix in
Section 4.3 for details. Models are evaluated on the test set (January
2024 -- December 2025, 24 months) using four metrics:

\begin{itemize}
\tightlist
\item
  \textbf{MAE} (Mean Absolute Error)
\item
  \textbf{RMSE} (Root Mean Squared Error)
\item
  \textbf{MAPE} (Mean Absolute Percentage Error)
\item
  \textbf{R\(^2\)} (Coefficient of Determination): Proportion of
  variance explained, computed as
  \(R^2 = 1 - \frac{\sum_i (y_i - \hat{y}_i)^2}{\sum_i (y_i - \bar{y}_{\mathrm{train}})^2}\),
  where \(\bar{y}_{\mathrm{train}}\) is the training-set mean. This
  follows standard time-series practice: the training-set mean serves as
  the naive baseline against which model predictions are evaluated.
  Negative values would indicate performance worse than predicting the
  training-set mean.
\end{itemize}

\subsection{Linear Regression}\label{linear-regression}

Linear Regression assumes a linear relationship between features
\(\mathbf{x}\) and target \(y\):

\[y = \mathbf{w}^T \mathbf{x} + b\]

The model minimizes the sum of squared residuals (Ordinary Least
Squares):

\[\min_{\mathbf{w}, b} \sum_{i=1}^{n} (y_i - \mathbf{w}^T \mathbf{x}_i - b)^2\]

\textbf{Assumptions:} Linearity, independence of residuals,
homoscedasticity, normally distributed errors. With 132 training samples
and 7 features (including intercept), the observation-to-parameter ratio
(132:8 \(\approx\) 16:1) is adequate.

\textbf{Applicability:} The \texttt{time\_idx} feature provides a strong
linear signal for the overall upward trend. With 132 training samples,
the model's parsimony is an advantage over complex alternatives
{[}11{]}.

\textbf{Limitations:} Cannot capture nonlinear patterns, seasonal cycles
beyond what month-encoded features provide, or structural breaks.

\subsection{Random Forest Regression}\label{random-forest-regression}

Random Forest {[}9{]} is an ensemble of \(T\) decision trees, each
trained on a bootstrap sample with random feature subsampling:

\[\hat{y} = \frac{1}{T} \sum_{t=1}^{T} h_t(\mathbf{x})\]

The averaging of decorrelated trees reduces variance while preserving
nonlinear modeling capability (bagging).

\textbf{Assumptions:} i.i.d. samples. Does not require stationarity or
linearity. \textbf{Key hyperparameters:} \texttt{n\_estimators} = 200,
\texttt{max\_depth} = 10 (to limit overfitting on 132 samples),
\texttt{min\_samples\_split} = 2.

\textbf{Advantage over quarterly analysis:} With 132 training samples
(vs.~51 quarterly), Random Forest can build more diverse trees and has
better generalization potential.

\subsection{XGBoost Regressor}\label{xgboost-regressor}

XGBoost {[}10{]} builds trees \textbf{sequentially}, with each tree
correcting errors of the previous ensemble:

\[\hat{y}_i = \sum_{k=1}^{K} f_k(\mathbf{x}_i)\]

At each step, a new tree minimizes a regularized objective:

\[\mathcal{L}^{(k)} = \sum_{i=1}^{n} \ell(y_i, \hat{y}_i^{(k-1)} + f_k(\mathbf{x}_i)) + \Omega(f_k)\]

where \(\ell\) is squared error and \(\Omega\) penalizes tree
complexity.

\textbf{Assumptions:} i.i.d. samples. Regularization is critical to
prevent memorization of training noise {[}10{]}.

\textbf{Key hyperparameters:} \texttt{n\_estimators} = 200,
\texttt{max\_depth} = 6, \texttt{learning\_rate} = 0.1.

\subsection{\texorpdfstring{SARIMAX\((1,1,1)(1,1,1)_{12}\)}{SARIMAX(1,1,1)(1,1,1)\_\{12\}}}\label{sarimax111111_12}

SARIMAX (Seasonal AutoRegressive Integrated Moving Average with
eXogenous variables) {[}7, 8{]} is a classical statistical model
decomposing a time series into autoregressive, differencing,
moving-average, and seasonal components, augmented with exogenous
regressors. To stabilize variance and ensure non-negative
back-transforms, the model is fitted to \(z_t = \log(y_t + 1)\):

\[\Phi_P(B^s)\,\phi_p(B)\,(1-B^s)^D\,(1-B)^d\, z_t = \sum_{j} \beta_j x_{j,t} + \Theta_Q(B^s)\,\theta_q(B)\,\varepsilon_t\]

where \(B\) is the backshift operator (\(B y_t = y_{t-1}\)), \(s = 12\)
is the seasonal period, \(x_{j,t}\) are exogenous variables, and
\(\varepsilon_t\) is white noise.

The model SARIMAX\((1,1,1)(1,1,1)_{12}\) was selected via AIC-minimizing
grid search over \((p,d,q)(P,D,Q)_{12}\) combinations. It applies
first-order differencing at both regular and seasonal levels, with one
autoregressive and one moving-average term at each level. A binary
\texttt{covid\_closed} exogenous variable (1 for April 2020--December
2021, 0 otherwise) signals the structural break caused by border
closures, allowing the model to learn the level shift without
discontinuity in the calendar. For the 2026 forecast, the model was
refitted on all data through December 2025 before generating
12-month-ahead predictions.

\textbf{Assumptions:} Stationarity after differencing; white-noise
residuals. The COVID closure period (April 2020--December 2021) is
zero-imputed with \texttt{covid\_closed\ =\ 1}, preserving calendar
continuity so that lag features bridge the gap correctly.

\textbf{Advantage of monthly data:} With \(s = 12\), the model captures
finer seasonal patterns (e.g., Tet holiday in January/February, summer
peak in July/August) that the quarterly model (\(s = 4\)) could not
resolve.

\subsection{Chronos-T5}\label{chronos-t5}

Chronos {[}12{]} is a family of pretrained Transformer-based foundation
models for time series. Unlike the other models, Chronos does not learn
from the 132 training samples; it was pretrained on millions of time
series from diverse domains. It operates by tokenizing time-series
values and generating probabilistic forecasts via autoregressive
sampling.

\begin{longtable}[]{@{}lll@{}}
\toprule\noalign{}
Model & Parameters & Context window \\
\midrule\noalign{}
\endhead
\bottomrule\noalign{}
\endlastfoot
chronos-t5-tiny & 8M & 512 tokens \\
chronos-t5-small & 46M & 512 tokens \\
chronos-t5-base & 200M & 512 tokens \\
\end{longtable}

\textbf{Advantage of monthly data:} With 132 context points (vs.~55
quarterly), the model has a denser signal for pattern recognition. The
512-token context window is well within limits.

\subsection{CIR\# Stochastic Differential Equation
Model}\label{cir-stochastic-differential-equation-model}

The CIR\# model {[}13, 14{]} extends the Cox--Ingersoll--Ross SDE for
tourism forecasting with disrupted data:

\[dr(t) = \kappa(\theta - r(t))\,dt + \sigma\sqrt{r(t)}\,dW(t)\]

where \(r(t)\) denotes monthly aggregate tourist arrivals expressed as a
count (persons), \(\kappa > 0\) is the mean-reversion speed (per month),
\(\theta\) is the long-run mean level (persons/month), \(\sigma\) is the
volatility parameter with units persons/month\(^{1/2}\), and \(dW(t)\)
is a standard Brownian motion increment. The scaling of \(r(t)\) is raw
monthly arrival counts (not log-transformed or rescaled). In discrete
form with \(\Delta t = 1\) month, the Euler approximation is
\(\Delta r_t = \kappa(\theta - r_t) + \sigma\sqrt{r_t}\,\epsilon_t\)
with \(\epsilon_t \sim \mathcal{N}(0,1)\). Parameters are estimated via
OLS of \(\Delta r_t\) on \(r_t\):

\[\Delta r_t = \alpha + \beta r_t + \epsilon_t\]

yielding \(\hat{\alpha} = \hat{\kappa}\hat{\theta}\),
\(\hat{\kappa} = -\hat{\beta}\), and
\(\hat{\theta} = \hat{\alpha}/\hat{\kappa}\) when \(\hat{\kappa} > 0\).
The volatility is estimated from the residual standard deviation:
\(\hat{\sigma} = \text{SD}(\hat{\epsilon}_t) / \sqrt{\bar{r}_t}\).

Orlando and Bufalo {[}13{]} report MAPE of 1.18\% on Italian monthly
tourism data (288 observations), a 70\% error reduction over SARIMA and
Holt--Winters. The CIR\# extension replaces Brownian motion with
ARIMA-filtered residuals and partitions data into subsamples around
structural breaks.

\textbf{Expected failure on Vietnam's data:} CIR\# was included
specifically to test its documented boundary condition {[}13, 14{]}: the
model assumes stationary, mean-reverting processes. The fitted
regression yields \(\hat{\beta} < 0\), implying \(\hat{\kappa} > 0\)
(apparent mean-reversion in-sample). However, Vietnam's tourism exhibits
a strong secular upward trend that violates the stationarity assumption:
the Augmented Dickey-Fuller test on the training series cannot reject a
unit root at the 5\% significance level (\(p = 0.079\)). This means the
apparent \(\hat{\kappa} > 0\) is a finite-sample artifact rather than
genuine mean-reversion. The model consequently produces forecasts that
revert toward the in-sample mean rather than following the post-COVID
recovery trajectory, confirming that CIR\# is poorly suited to this
strongly nonstationary setting.

\subsection{Model Comparison}\label{model-comparison}

\begin{figure}
\centering
\includegraphics{output/model_comparison.png}
\caption{Model performance comparison (MAE, MAPE, R\(^2\)).}
\end{figure}

\begin{longtable}[]{@{}lllll@{}}
\toprule\noalign{}
Model & MAE & RMSE & MAPE & R² \\
\midrule\noalign{}
\endhead
\bottomrule\noalign{}
\endlastfoot
Random Forest & 117,361 & 151,813 & 7.69\% & 0.963 \\
Linear Regression & 143,260 & 164,870 & 9.66\% & 0.956 \\
Chronos-T5-tiny & 172,046 & 221,354 & 11.16\% & 0.921 \\
XGBoost & 248,413 & 300,611 & 16.10\% & 0.854 \\
SARIMAX (log) & 331,011 & 370,909 & 21.52\% & 0.778 \\
CIR\# & 394,559 & 465,267 & 24.98\% & 0.650 \\
\end{longtable}

\emph{R² is computed relative to the training-set mean
(\(\bar{y}_{\mathrm{train}} = 717{,}881\)) following standard
time-series practice. All supervised ML models include Fourier seasonal
features (\(\sin(2\pi k \cdot \text{month}/12)\),
\(\cos(2\pi k \cdot \text{month}/12)\) for \(k=1,2\)) and a Tet holiday
month indicator. SARIMAX captures seasonality through its
\((1,1,1)_{12}\) seasonal order instead. Chronos uses raw time-series
values only.}

\textbf{Key observations:}

\begin{itemize}
\item
  \textbf{Random Forest achieves the lowest error} by both MAE (117,361)
  and MAPE (7.69\%). The Fourier terms and Tet calendar feature improved
  RF from 8.04\% to 7.69\% MAPE, with the largest gain for XGBoost
  (18.39\% → 16.10\%). Linear Regression achieves the highest R² (0.956)
  and competitive MAPE (9.66\%).
\item
  \textbf{Chronos-T5-tiny} (8.4M parameters) achieves a competitive MAPE
  of 11.16\% using only raw time-series values (no engineered features),
  demonstrating the power of transfer learning for tourism forecasting.
  A scaling analysis across tiny/small/base variants found that the
  smallest model performs best on this dataset: tiny (MAPE 11.16\%),
  base (13.20\%), small (13.89\%). This counterintuitive result is
  consistent with the small training window (132 months) being
  insufficient to benefit from larger model capacity.
\item
  \textbf{XGBoost underperforms} Random Forest (MAPE 16.10\% vs 7.69\%),
  consistent with gradient boosting's known sensitivity to small sample
  sizes {[}10, 11{]}. A learning curve analysis confirms non-monotonic
  behavior: MAPE varied from 17.6\% (50\% data) to 34.2\% (75\%) to
  18.4\% (100\%), indicating high variance and instability.
\item
  \textbf{Feature importance} (Random Forest): \texttt{lag\_1}
  dominates, followed by \texttt{lag\_12} and \texttt{time\_idx}. The
  Fourier terms rank below the lag features but above \texttt{year} in
  importance.
\item
  \textbf{Multi-step evaluation results.} A rolling-origin evaluation
  with horizons \(h = 1, 3, 6, 12\) confirms that Random Forest
  maintains remarkably stable accuracy across horizons (MAPE \(\approx\)
  9--10\%), while Linear Regression and XGBoost degrade significantly at
  longer horizons (MAPE \(\approx\) 30--31\% at \(h = 12\)). This is
  because RF's bagging approach provides robustness against error
  accumulation, whereas LR's linear extrapolation and XGBoost's
  sequential boosting amplify forecast drift.
\item
  \textbf{SARIMAX} (MAPE = 21.52\%) with log-transformed target
  \(\log(y_t+1)\) and \texttt{covid\_closed} exogenous variable. The
  model was selected via grid search over \((p,d,q)(P,D,Q)_{12}\) orders
  minimizing AIC, yielding SARIMAX\((1,1,1)(1,1,1)_{12}\). The
  log-transformation stabilizes variance and ensures non-negative
  back-transforms, but the model still struggles to extrapolate the
  post-COVID growth trend from a stationary-process framework.
\item
  \textbf{CIR\#} (MAPE = 24.98\%) was included to test its boundary
  condition. The estimated \(\kappa\) is positive, but the strongly
  trending data violates the stationarity assumption (\(p = 0.079\) on
  the ADF test), producing mean-reverting forecasts that fail to capture
  post-COVID growth.
\end{itemize}

\newpage

\section{Model Optimization and External
Features}\label{model-optimization-and-external-features}

\subsection{Hyperparameter
Optimization}\label{hyperparameter-optimization}

\textbf{Random Forest} was optimized via GridSearchCV with
\texttt{TimeSeriesSplit(n\_splits=3)} to prevent temporal leakage:

\begin{itemize}
\tightlist
\item
  \texttt{n\_estimators}: \{100, 200, 300\}
\item
  \texttt{max\_depth}: \{5, 10, 15, None\}
\item
  \texttt{min\_samples\_split}: \{2, 5, 10\}
\end{itemize}

\textbf{XGBoost} was optimized via RandomizedSearchCV (50 iterations)
with \texttt{TimeSeriesSplit(n\_splits=3)}:

\begin{itemize}
\tightlist
\item
  \texttt{n\_estimators}: 100--500
\item
  \texttt{max\_depth}: 3--9
\item
  \texttt{learning\_rate}: 0.01--0.2
\item
  \texttt{subsample}: 0.7--1.0
\item
  \texttt{colsample\_bytree}: 0.7--1.0
\end{itemize}

\subsection{External Features}\label{external-features}

Exchange rates (VND vs.~KRW, CNY, USD, JPY, TWD, MYR, THB, RUB) were
obtained from Yahoo Finance {[}17{]} as end-of-month spot rates. Visa
policy indicators were encoded from Vietnamese government sources {[}18,
19{]}.

\textbf{Result:} External features provided marginal improvement for
Linear Regression but degraded tree-based model performance due to
overfitting. Feature importance analysis shows that \texttt{lag\_1} and
\texttt{rolling\_mean\_12} remain the dominant features. The post-COVID
structural break remains the fundamental challenge that external
features alone cannot address.

\subsection{Fourier Seasonal Features and Tet Calendar
Effect}\label{fourier-seasonal-features-and-tet-calendar-effect}

To explicitly model seasonality beyond what the \texttt{month} integer
feature provides, we added Fourier terms (two harmonics, \(k=1,2\)):

\[f_k^{(\sin)}(t) = \sin\!\left(\frac{2\pi k \cdot m_t}{12}\right), \qquad f_k^{(\cos)}(t) = \cos\!\left(\frac{2\pi k \cdot m_t}{12}\right)\]

where \(m_t \in \{1,\ldots,12\}\) is the calendar month. These four
features capture the dominant seasonal cycle and its first harmonic.

Additionally, a binary \texttt{tet\_month} indicator was added to
capture the Tet holiday (Lunar New Year) effect, which shifts between
January and February depending on the lunar calendar. Tet falls in
January for years 2012, 2014, 2017, 2020, 2023 and in February for years
2013, 2015, 2016, 2018, 2019, 2021, 2022, 2024, 2025, 2026.

\textbf{Impact on model accuracy:}

\begin{longtable}[]{@{}llll@{}}
\toprule\noalign{}
Model & Base MAPE & +Fourier+Tet MAPE & Improvement \\
\midrule\noalign{}
\endhead
\bottomrule\noalign{}
\endlastfoot
Linear Regression & 9.73\% & 9.66\% & \$-\$0.07 pp \\
Random Forest & 8.04\% & 7.69\% & \$-\$0.35 pp \\
XGBoost & 18.39\% & 16.10\% & \$-\$2.29 pp \\
\end{longtable}

The Fourier terms and Tet indicator provided the largest improvement for
XGBoost (2.29 percentage points), suggesting that explicit seasonal
features help gradient boosting capture patterns that would otherwise
require deeper trees. SARIMAX captures seasonality through its
\((1,1,1)_{12}\) seasonal order and does not benefit from Fourier terms.

\subsection{Chronos Scaling Behavior
Analysis}\label{chronos-scaling-behavior-analysis}

We evaluated three Chronos-T5 variants of increasing size on the same
test set to characterize scaling behavior:

\begin{longtable}[]{@{}llll@{}}
\toprule\noalign{}
Variant & Parameters & MAPE & R² \\
\midrule\noalign{}
\endhead
\bottomrule\noalign{}
\endlastfoot
chronos-t5-tiny & 8.4M & 11.16\% & 0.921 \\
chronos-t5-small & 46.2M & 13.89\% & 0.880 \\
chronos-t5-base & 201.4M & 13.20\% & 0.926 \\
\end{longtable}

The smallest model (tiny) achieves the best MAPE, while the largest
model (base) achieves the best R². This counterintuitive result---larger
models do not systematically improve accuracy---is consistent with the
limited training context (132 months): the 512-token context window of
all three variants can accommodate the full training history, so the
additional parameters in larger models do not provide more information.
The base model's slightly higher R² suggests better calibration of
probabilistic forecasts, but its MAPE is 2 percentage points worse than
tiny. For operational deployment on this dataset, chronos-t5-tiny offers
the best accuracy-efficiency tradeoff.

\subsection{XGBoost Learning Curve
Analysis}\label{xgboost-learning-curve-analysis}

To investigate XGBoost's underperformance, we evaluated MAPE at
different training sample fractions:

\begin{longtable}[]{@{}lll@{}}
\toprule\noalign{}
Training fraction & N samples & MAPE \\
\midrule\noalign{}
\endhead
\bottomrule\noalign{}
\endlastfoot
25\% & 33 & 58.01\% \\
50\% & 66 & 17.61\% \\
75\% & 99 & 34.15\% \\
100\% & 132 & 18.39\% \\
\end{longtable}

The non-monotonic behavior (75\% data yields worse MAPE than 50\%)
confirms that XGBoost exhibits high variance on this dataset. This is
characteristic of gradient boosting with limited training data: the
sequential error-correction mechanism amplifies noise when sample sizes
are insufficient to reliably estimate residuals. The learning curve does
not show clear convergence, suggesting that XGBoost would require
substantially more training data to match Random Forest's performance.

\newpage

\section{Forecasting and Evaluation}\label{forecasting-and-evaluation}

\subsection{Predicted vs.~Actual (Test
Set)}\label{predicted-vs.-actual-test-set}

\begin{figure}
\centering
\includegraphics{output/pred_vs_actual.png}
\caption{Predicted vs.~actual monthly arrivals for the test set
(2024--2025).}
\end{figure}

\textbf{Forecasting methodology note.} All four models generate
12-month-ahead forecasts for 2026. The tree-based models (Linear
Regression, Random Forest, XGBoost) use a recursive strategy: each
month's prediction is fed back as the \texttt{lag\_1} feature for the
next month. This accumulates error at each step but captures the
nonlinear patterns the tree models learned. SARIMAX uses its
autoregressive structure to generate multi-step predictions directly.
The ensemble mean (shown in red) averages all four models, providing a
more robust forecast than any single model. The shaded band shows the
range between the most optimistic and most pessimistic single-model
predictions---a model-disagreement band showing inter-model forecast
uncertainty. SARIMAX log-scale intervals are reported separately and
widen rapidly over the 12-month horizon due to the multiplicative nature
of the log-transform; these are conceptually distinct from the ensemble
disagreement band.

\subsection{SARIMAX 12-Month Forecast
(2026)}\label{sarimax-12-month-forecast-2026}

\begin{figure}
\centering
\includegraphics{output/forecast_plot.png}
\caption{SARIMAX 12-month forecast for 2026 with 95\% confidence
interval.}
\end{figure}

\emph{Figure shows the SARIMAX-only forecast (refitted through December
2025) with 95\% confidence intervals. The ensemble point forecast is
reported separately below.}

\begin{longtable}[]{@{}llll@{}}
\toprule\noalign{}
Month & Forecast & 95\% CI Lower & 95\% CI Upper \\
\midrule\noalign{}
\endhead
\bottomrule\noalign{}
\endlastfoot
Jan 2026 & 1,769,400 & 1,065,283 & 2,938,916 \\
Feb 2026 & 1,851,450 & 803,824 & 4,264,449 \\
Mar 2026 & 1,647,989 & 542,942 & 5,002,129 \\
Apr 2026 & 1,540,470 & 400,224 & 5,929,291 \\
May 2026 & 1,499,244 & 316,193 & 7,108,713 \\
Jun 2026 & 1,479,019 & 258,836 & 8,451,254 \\
Jul 2026 & 1,609,229 & 237,686 & 10,895,052 \\
Aug 2026 & 1,827,948 & 230,939 & 14,468,645 \\
Sep 2026 & 1,676,192 & 183,104 & 15,344,327 \\
Oct 2026 & 1,729,753 & 164,843 & 18,150,715 \\
Nov 2026 & 1,910,191 & 160,010 & 22,803,647 \\
Dec 2026 & 1,994,000 & 147,774 & 26,905,886 \\
\end{longtable}

\textbf{Ensemble point forecast.} The ensemble forecast averages the
recursive multi-step predictions of Linear Regression, Random Forest,
XGBoost, and the SARIMAX median back-transform. Full ensemble table
available in \texttt{output/ensemble\_forecast.csv}.

The back-transform \(\exp(\hat{z}_t) - 1\) produces a
\emph{median-scale} forecast on the original scale, not the conditional
mean. An unbiased mean forecast would require a lognormal correction
\(\exp(\hat{z}_t + \hat{\sigma}^2_t/2) - 1\) or simulation-based
averaging. The confidence intervals are asymmetric on the original
scale, reflecting the multiplicative nature of the log-transform. The
widening of the intervals (e.g., December 2026 upper bound exceeding
26M) indicates that model uncertainty grows rapidly over the 12-month
horizon, consistent with the challenges of multi-step forecasting with
log-transformed targets.

\subsection{Per-Country Forecasts
(2026)}\label{per-country-forecasts-2026}

Ensemble models (averaging Linear Regression, Random Forest, XGBoost,
and SARIMAX) were fitted to the top 5 source countries individually. The
tree-based models use recursive forecasting to generate 12-month-ahead
predictions; SARIMAX uses its autoregressive structure directly.

\begin{figure}
\centering
\includegraphics{output/country_forecasts_plot.png}
\caption{Ensemble forecasts for top 5 source countries (2026). Model
range band shows disagreement between LR, RF, XGBoost, and SARIMAX.}
\end{figure}

\begin{longtable}[]{@{}llllll@{}}
\toprule\noalign{}
Month & Hàn Quốc & Trung Quốc & Đài Loan & Nhật Bản & Hoa Kỳ \\
\midrule\noalign{}
\endhead
\bottomrule\noalign{}
\endlastfoot
Jan 2026 & 415,589 & 253,357 & 90,615 & 64,534 & 76,411 \\
Feb 2026 & 415,133 & 245,562 & 102,170 & 63,645 & 70,282 \\
Mar 2026 & 364,367 & 213,994 & 87,645 & 67,327 & 59,654 \\
Apr 2026 & 365,820 & 223,717 & 92,212 & 62,500 & 58,413 \\
May 2026 & 362,484 & 221,843 & 91,027 & 64,439 & 52,122 \\
Jun 2026 & 374,440 & 212,771 & 93,809 & 60,466 & 64,743 \\
Jul 2026 & 380,907 & 225,928 & 102,468 & 64,545 & 68,922 \\
Aug 2026 & 434,960 & 246,152 & 99,805 & 80,036 & 61,241 \\
Sep 2026 & 396,678 & 238,087 & 97,347 & 79,496 & 55,602 \\
Oct 2026 & 404,532 & 239,295 & 101,391 & 71,685 & 58,810 \\
Nov 2026 & 406,612 & 252,036 & 102,543 & 75,630 & 65,001 \\
Dec 2026 & 438,882 & 249,132 & 103,948 & 77,370 & 65,566 \\
\end{longtable}

\emph{Table shows ensemble mean of Linear Regression, Random Forest,
XGBoost, and SARIMAX. Model disagreement range: 8.2M--10.1M total.}

The top 5 countries account for 73.7\% of the total 2026 ensemble
forecast (10.3M of 14.0M). Hàn Quốc (South Korea) is projected to remain
the largest source market, followed by Trung Quốc. The shaded band shows
the range between the highest and lowest model predictions --- a more
interpretable measure of uncertainty than the SARIMAX confidence
interval.

\newpage

\subsection{Forecast Validation (Jan--May
2026)}\label{forecast-validation-janmay-2026}

With actual data available for the first five months of 2026, we
evaluate the SARIMAX-only forecast (refitted through December 2025)
against the 32-country source-country sums (\(Y^{32}_t\)) from
\texttt{df\_monthly.csv}. The MAPE formula used is:

\[\text{MAPE} = \frac{1}{n}\sum_{i=1}^{n} \left|\frac{y_i - \hat{y}_i}{y_i}\right| \times 100\%\]

\begin{figure}
\centering
\includegraphics{output/forecast_validation.png}
\caption{Forecast vs Actual for Jan---May 2026 (aggregate and top source
countries).}
\end{figure}

\begin{longtable}[]{@{}lllll@{}}
\toprule\noalign{}
Month & Actual (\(Y^{32}_t\)) & SARIMAX Forecast & Error & Abs \%
Error \\
\midrule\noalign{}
\endhead
\bottomrule\noalign{}
\endlastfoot
Jan 2026 & 2,210,922 & 1,769,400 & \$-\$20.0\% & 20.0\% \\
Feb 2026 & 1,970,499 & 1,851,450 & \$-\$6.0\% & 6.0\% \\
Mar 2026 & 1,877,290 & 1,647,989 & \$-\$12.2\% & 12.2\% \\
Apr 2026 & 1,828,573 & 1,540,470 & \$-\$15.8\% & 15.8\% \\
May 2026 & 1,606,802 & 1,499,244 & \$-\$6.7\% & 6.7\% \\
\textbf{MAPE} & & & & \textbf{12.1\%} \\
\end{longtable}

\emph{Actuals are \(Y^{32}_t\) (32-country sums from
\texttt{df\_monthly.csv}). Forecast source:
SARIMAX(1,1,1)(1,1,1)\(_{12}\) refitted on training data 2013--2025,
then forecast 12 steps ahead. Note: official GSO all-market totals
(\(Y^{\mathrm{official}}_t\)) include residual categories not in the
32-country set and would be higher.}

\textbf{Per-country validation:}

\begin{longtable}[]{@{}
  >{\raggedright\arraybackslash}p{(\columnwidth - 4\tabcolsep) * \real{0.4091}}
  >{\raggedright\arraybackslash}p{(\columnwidth - 4\tabcolsep) * \real{0.2727}}
  >{\raggedright\arraybackslash}p{(\columnwidth - 4\tabcolsep) * \real{0.3182}}@{}}
\toprule\noalign{}
\begin{minipage}[b]{\linewidth}\raggedright
Country
\end{minipage} & \begin{minipage}[b]{\linewidth}\raggedright
MAPE
\end{minipage} & \begin{minipage}[b]{\linewidth}\raggedright
Notes
\end{minipage} \\
\midrule\noalign{}
\endhead
\bottomrule\noalign{}
\endlastfoot
Hàn Quốc (South Korea) & 8.3\% & Best fit; model captures seasonal
pattern well \\
Trung Quốc & 48.1\% & Consistent underprediction; structural growth
since 2024 not captured \\
Nhật Bản & 25.9\% & Improved with corrected parser; seasonal pattern
partially captured \\
Campuchia & 50.3\% & Land-border regime shift: visa exemptions + new air
routes + healthcare tourism; Jan 2026 hit 223,000 (3\(\times\) Dec
2025) \\
\end{longtable}

\textbf{Russia case study.} The SARIMAX component forecasts
7,000--12,000 arrivals/month for Russia in 2026 (the ensemble mean in
the table is higher because the other three models pull the average up);
actual figures are 113,000--137,000 --- a 10\(\times\) error. This is an
out-of-distribution forecast failure caused by a geopolitical regime
shift that was not represented in the 2013--2023 training data. The
Russia-Ukraine war (2022) closed EU and US destinations to Russian
tourists. Vietnam, with its 45-day visa exemption, resumed direct
flights (Aeroflot, VietJet, Vietnam Airlines), and affordable pricing,
became a primary alternative. Russian arrivals grew from 39,921 (2022,
per the GSO monthly dataset; the VNAT-published annual total for 2022 is
28,056) to 689,714 (2025), surpassing the pre-pandemic peak of 646,524
(2019). No model trained on pre-2024 data can anticipate this structural
shift.

\textbf{Campuchia case study.} The model forecasts 27,000--38,000
Cambodian arrivals/month for 2026; actual figures range from 53,000 to
223,000. Two effects compound:

\begin{enumerate}
\def\labelenumi{\arabic{enumi}.}
\item
  \textbf{Baseline regime shift.} Cambodia's monthly average rose from
  \textasciitilde38,000 (2024) to \textasciitilde57,000 (2025--2026),
  driven by reciprocal 30-day visa exemptions, Air Cambodia's new routes
  (\textasciitilde10 daily flights to Vietnam), and accelerating
  cross-border healthcare tourism and shopping. Cambodia's proximity
  (6-hour bus from Phnom Penh to Ho Chi Minh City) makes it
  fundamentally different from air-dependent markets.
\item
  \textbf{Tet seasonal amplification.} January spikes recur: January
  2025 hit 100,000 (+64\% over December); January 2026 hit 223,025
  (+206\% over December). Cambodians cross the border en masse for Lunar
  New Year shopping and family visits. The 2026 spike was double the
  2025 spike due to new flight capacity launched in late 2025.
  February--May 2026 settled to 53,000--66,000/month, confirming the
  elevated baseline.
\end{enumerate}

The model, trained on 2013--2023 data showing Cambodia at
20,000--70,000/month, cannot anticipate either effect {[}21{]}.

The five-month MAPE of 12.1\% shows that the corrected SARIMAX model
(refitted through December 2025) achieves substantially better accuracy
than the previous estimate of 34.8\% (which used an uncorrected model
fitted only through December 2023). Systematic underprediction persists,
however, reflecting the post-COVID growth acceleration not fully
captured by the stationary-process framework. Hàn Quốc (South Korea) is
the best-predicted country (MAPE = 8.3\%), as its growth trajectory most
closely resembles the training distribution.

\section{Conclusion}\label{conclusion}

\subsection{Key Findings}\label{key-findings}

\begin{enumerate}
\def\labelenumi{\arabic{enumi}.}
\item
  \textbf{Monthly data significantly improves analysis.} With 132
  training observations (vs.~51 quarterly), machine learning models have
  more than double the data for learning patterns. SARIMAX with
  \(s = 12\) resolves finer seasonal structure (Tet holiday, summer
  peaks) that \(s = 4\) could not capture.
\item
  \textbf{Coverage bias persists.} Only 11--13 countries reported
  monthly data during 2009--2011, compared to 29--31 from December 2012
  onward. Aggregate trend analyses should use January 2013 as the
  starting point to ensure stable country coverage throughout the
  training period.
\item
  \textbf{China and South Korea dominate} the tourism market, accounting
  for the largest cumulative arrivals. Emerging markets (Hong Kong,
  Spain, Italy, Philippines) grew fastest during 2013--2019.
\item
  \textbf{Strong seasonality:} January--February consistently has the
  highest arrivals (Tet holiday + winter tourism), with a secondary
  summer peak in June--August.
\item
  \textbf{Post-COVID structural break} remains the dominant challenge.
  When R² is computed relative to the training-set mean (standard
  time-series practice), all models achieve positive R² because
  predictions capture the upward trend from the training level
  (\(\bar{y}_{\mathrm{train}} = 717{,}881\)) to the much higher test-set
  range. The corrected model comparison shows Random Forest achieves the
  best MAPE (7.69\%) and MAE (117,361), followed by Linear Regression
  (MAPE 9.66\%, R² = 0.956). Chronos-T5-tiny achieves competitive MAPE
  (11.16\%) with no task-specific training, and a scaling analysis found
  the smallest variant outperforms both small and base on this 132-month
  dataset. XGBoost (MAPE 16.10\%), SARIMAX (MAPE 21.52\%), and CIR\#
  (MAPE 24.98\%) underperform. The test distribution (2024--2025
  recovery) differs fundamentally from the 2013--2023 training
  distribution, penalizing models that cannot extrapolate the growth
  trend.
\item
  \textbf{CIR\# fails on trending data.} Despite having monthly data as
  recommended by Orlando and Bufalo {[}13{]}, the model's mean-reversion
  assumption is violated by Vietnam's upward-trending tourism trajectory
  (ADF test: \(p = 0.079\)). This constitutes a documented boundary
  condition for the model {[}13, 14{]}.
\item
  \textbf{Fourier terms and Tet calendar features provide measurable
  improvement.} Adding \(\sin(2\pi k \cdot m/12)\) and
  \(\cos(2\pi k \cdot m/12)\) (\(k=1,2\)) and a binary Tet holiday month
  indicator reduced MAPE by 0.07--2.29 percentage points across the
  supervised models, with the largest gain for XGBoost (18.39\% →
  16.10\%). Exchange rates and visa policy indicators provided only
  marginal additional improvement.
\end{enumerate}

\subsection{Limitations}\label{limitations}

\begin{enumerate}
\def\labelenumi{\arabic{enumi}.}
\tightlist
\item
  \textbf{COVID gap (2020--2021):} Zero-imputation of the COVID closure
  period preserves calendar continuity; the \texttt{covid\_closed}
  exogenous variable signals the structural break to SARIMAX. However,
  the zero-filled months still represent a departure from the true
  data-generating process. Sensitivity analysis excluding these months
  or using alternative imputation (e.g., last non-COVID value) would
  strengthen robustness claims.
\item
  \textbf{Limited training set (132 months):} Starting from January 2013
  for stable country coverage reduces the training window from 144 to
  132 months. While substantially better than 51 quarterly points, this
  still constrains model complexity compared to the 288 monthly
  observations used in the Italian CIR\# study {[}13{]}.
\item
  \textbf{Missing external features:} Source-country GDP, flight
  capacity, Google Trends search volume, and oil prices were not
  included. These may improve post-COVID extrapolation accuracy.
\item
  \textbf{Coverage inconsistency (2009--2011):} Only 11--13 countries
  per month limits reliability of aggregate statistics for early years.
  Pre-entry missingness is treated as zero arrivals, which may introduce
  artificial growth jumps when countries enter the reporting framework.
\item
  \textbf{Formal diagnostics performed:} ADF test on the training series
  yields \(p = 0.079\) (borderline non-stationary at 5\%); KPSS test
  yields \(p = 0.076\) (borderline stationary at 5\%); the differenced
  series is clearly stationary (\(p < 0.001\)). Ljung-Box test on
  SARIMAX residuals shows no autocorrelation at lag 6 (\(p = 0.62\)) but
  significant seasonal autocorrelation at lag 12 (\(p = 0.002\)).
  Breusch-Pagan test confirms heteroscedasticity (\(p = 0.002\)).
  Jarque-Bera test confirms non-normal residuals (\(p < 0.001\)). These
  findings suggest the SARIMAX model could be improved with additional
  seasonal terms or heteroscedasticity-robust standard errors.
\item
  \textbf{Recursive multi-step forecasting limitations:} The 2026
  forecasts use recursive prediction, where each month's forecast feeds
  back as the \texttt{lag\_1} feature for the next month. Rolling-origin
  evaluation shows error accumulation at longer horizons: Random Forest
  remains stable (MAPE \textasciitilde10\% at \(h = 12\)) but Linear
  Regression and XGBoost degrade to MAPE \textasciitilde30\% at
  \(h = 12\).
\item
  \textbf{No ensembling strategies beyond averaging:} The ensemble
  forecast uses simple arithmetic mean. Stacked generalization, weighted
  averaging (based on recent performance), or model selection strategies
  could improve accuracy.
\end{enumerate}

\subsection{Future Directions}\label{future-directions}

\begin{enumerate}
\def\labelenumi{\arabic{enumi}.}
\tightlist
\item
  \textbf{Source-country GDP and flight capacity} as leading indicators
\item
  \textbf{Google Trends search volume} for destination queries
\item
  \textbf{Regime-switching models} that explicitly handle pre-COVID,
  during-COVID, and post-COVID transitions
\item
  \textbf{Per-country models} (different countries exhibit different
  recovery trajectories)
\item
  \textbf{Ensemble methods} combining Linear Regression (trend) with
  Chronos (pattern recognition)
\item
  \textbf{Formal residual diagnostics} and stationarity testing
\end{enumerate}

\newpage

\section{References}\label{references}

{[}1{]} World Travel and Tourism Council, ``Economic Impact Research:
Vietnam,'' WTTC, 2024. {[}Online{]}. Available:
https://wttc.org/research/economic-impact

{[}2{]} Vietnam National Authority of Tourism, ``Báo cáo tình hình kinh
tế -- xã hội năm 2019: Khách quốc tế đến Việt Nam,'' VNAT, Hanoi, 2020.
{[}Online{]}. Available: https://vietnamtourism.gov.vn/

{[}3{]} H. T. Thi, ``Vietnam Tourism Industry During COVID-19
Pandemic,'' 2022. See also: General Statistics Office of Vietnam,
``Socio-economic situation report 2020,'' GSO, Hanoi, 2021.

{[}4{]} General Statistics Office of Vietnam, ``Tình hình kinh tế -- xã
hội năm 2025: Khách quốc tế đạt 21,2 triệu lượt,'' GSO, Hanoi,
Jan.~2026. {[}Online{]}. Available: https://www.gso.gov.vn/

{[}5{]} H. Song and S. F. Witt, \emph{Tourism Demand Modelling and
Forecasting: Modern Econometric Approaches}. Routledge, 2000.

{[}6{]} H. Song, G. Li, and Z. Cao, ``Tourism demand modelling and
forecasting --- A review of recent research,'' \emph{Tourism
Management}, vol.~74, pp.~217--232, 2019.

{[}7{]} R. J. Hyndman and G. Athanasopoulos, \emph{Forecasting:
Principles and Practice}, 3rd ed.~OTexts, 2021. {[}Online{]}. Available:
https://otexts.com/fpp3/

{[}8{]} G. E. P. Box, G. M. Jenkins, G. C. Reinsel, and G. M. Ljung,
\emph{Time Series Analysis: Forecasting and Control}, 5th ed.~Wiley,
2016.

{[}9{]} L. Breiman, ``Random Forests,'' \emph{Machine Learning},
vol.~45, no. 1, pp.~5--32, 2001.

{[}10{]} T. Chen and C. Guestrin, ``XGBoost: A Scalable Tree Boosting
System,'' in \emph{Proc. 22nd ACM SIGKDD Int. Conf. Knowledge Discovery
and Data Mining}, 2016, pp.~785--794.

{[}11{]} T. Hastie, R. Tibshirani, and J. Friedman, \emph{The Elements
of Statistical Learning}, 2nd ed.~Springer, 2009.

{[}12{]} A. F. Ansari et al., ``Chronos: Learning the Language of Time
Series,'' \emph{Trans. Mach. Learn. Res.}, 2024. {[}Online{]}.
Available: https://arxiv.org/abs/2403.07815

{[}13{]} G. Orlando and M. Bufalo, ``Improved tourism demand forecasting
with CIR\# model: the case of disrupted data patterns in Italy,''
\emph{Tourism Review}, vol.~79, no. 2, pp.~445--470, 2023, doi:
10.1108/TR-04-2023-0194.

{[}14{]} G. Orlando and M. Bufalo, ``The CIR\# model for time series
forecasting,'' \emph{Technological and Economic Development of Economy},
vol.~29, no. 5, pp.~1403--1427, 2023.

{[}15{]} General Statistics Office of Vietnam, ``Quarterly and monthly
international arrival statistics by country,'' GSO, Hanoi, 2008--May
2026 (partial). {[}Online{]}. Available: https://www.gso.gov.vn/

{[}16{]} R. J. A. Little and D. B. Rubin, \emph{Statistical Analysis
with Missing Data}, 2nd ed.~Wiley, 2002.

{[}17{]} Yahoo Finance, ``Historical exchange rates,'' Yahoo, 2025.
{[}Online{]}. Available: https://finance.yahoo.com/

{[}18{]} Vietnam Immigration Department, ``Pilot e-visa system for
foreign visitors,'' effective Feb.~1, 2017, initially for 40 countries.

{[}19{]} Vietnam National Assembly, ``Resolution on extension and
amendment of e-visa policy,'' approved Jun.~24, 2023, effective Aug.~15,
2023.

{[}20{]} General Statistics Office of Vietnam, ``Monthly international
arrivals from Russia 2022--2025,'' GSO, Hanoi, 2025. {[}Online{]}.
Available: https://www.gso.gov.vn/

{[}21{]} General Statistics Office of Vietnam, ``Monthly international
arrivals from Cambodia 2024--2026,'' GSO, Hanoi, 2026. See also:
VnExpress International, ``Cambodia replaces Taiwan to become Vietnam's
3rd largest source of tourists,'' Feb.~2026.

{[}22{]} Scikit-learn developers, ``Scikit-learn: Machine Learning in
Python,'' 2025. {[}Online{]}. Available: https://scikit-learn.org/

{[}23{]} T. Chen et al., ``XGBoost Documentation,'' 2025. {[}Online{]}.
Available: https://xgboost.readthedocs.io/

{[}24{]} Statsmodels developers, ``Statsmodels: Statistical Modeling and
Econometrics in Python,'' 2025. {[}Online{]}. Available:
https://www.statsmodels.org/

{[}25{]} Amazon Science, ``Chronos Forecasting,'' GitHub, 2024.
{[}Online{]}. Available:
https://github.com/amazon-science/chronos-forecasting

{[}26{]} Pacific Asia Travel Association, ``Visitor Arrivals to Asia
Pacific Destinations,'' PATA, H1 2019.

\end{document}
